
\subsection{Визуализация: Python-скрипты}

Скрипты находятся в \texttt{scripts/} и делятся на ядро (\texttt{core/})
и скрипты отрисовки (\texttt{visualization/}). Ядро содержит три компонента: \texttt{config.py} определяет все параметры отображения: цвета, размеры, колормапы. Неизменяемые
датаклассы гарантируют, что скрипт случайно не перезапишет глобальные
настройки. \texttt{GraphLoader} читает текстовые файлы в граф \textit{networkx}. \texttt{Renderer}~--- обёртка над
matplotlib/networkx с методами \texttt{setup\_plot}, \\\texttt{compute\_layout},
\texttt{draw\_nodes}, \texttt{draw\_edges}, \texttt{draw\_labels},
\texttt{finalize}. Это убирает дублирование конфигурации в каждом скрипте.

Каждый скрипт в \texttt{visualization/} принимает аргументы через
\texttt{argparse}, читает данные через \texttt{GraphLoader}, рисует через
\texttt{Renderer} и сохраняет в файл. \texttt{plot\_network.py}
обрабатывает оба типа~--- граф и сеть потоков~--- через флаг \texttt{--type}.
Толщина рёбра для графа пропорциональна весу, для потока~--- текущему потоку;
цвет ребра потока задаётся через колормап по коэффициенту загрузки
$flow/capacity$. \texttt{plot\_paths.py} получает файл с путями,
строит словарь рёбер, входящих в пути, и выделяет их цветом: каждый путь
своим цветом из \texttt{cm.rainbow}. \texttt{plot\_matrix.py} рисует
тепловую карту через seaborn с аннотациями значений.

\subsection{Веб-интерфейс}

Веб-интерфейс реализован на Streamlit (\texttt{web/}). Приложение
(\texttt{app.py}) запускает исполняемый файл C++ как дочерний процесс
через \texttt{subprocess.Popen} с открытым \texttt{stdin}. Взаимодействие
с C++ происходит через \texttt{send\_command}~--- метод просто пишет
строку в \texttt{proc.stdin}. C++ читает её как ввод пользователя в
консольном цикле и выполняет нужное действие.

Конфигурация описана в \texttt{config.py} через датаклассы.
\texttt{LabConfig} описывает лабораторную, \texttt{ActionConfig}~--- одно
действие внутри неё: номер команды (\texttt{execute\_cmd}), список
параметров ввода (\texttt{params}) и список файлов для отображения
(\texttt{images}). Класс \texttt{GraphLabApp} в \texttt{ui.py} обходит
эту конфигурацию и рендерит виджеты Streamlit динамически.

Дополнительно для лабораторной~1 строится гистограмма весов с наложенной
теоретической PDF распределения Рэлея--Райса. PDF вычисляется напрямую в
Python через \texttt{scipy.special.i0} (функция Бесселя) по той же формуле
Вадзинского, что и в C++. Это позволяет визуально сравнить реальное
распределение сгенерированных весов с теоретическим.
